\chapter{Images, render textures and colour}

Chapter 7 covered enough of rtextures to draw a sprite. This is the rest of the
module: editing pixels, drawing into textures, and the colour helpers.

\section{Image work happens on the CPU}

An \lstinline|Image| is pixels in ordinary RAM, and raylib gives you a small
image editor's worth of functions to work on them:

\begin{center}
\begin{tabular}{@{}ll@{}}
\toprule
\lstinline|ImageResize / ImageResizeNN| & smooth, or nearest-neighbour for pixel art \\
\lstinline|ImageCrop / ImageAlphaCrop| & cut a region, or trim transparent edges \\
\lstinline|ImageFlipHorizontal / Vertical| & mirror \\
\lstinline|ImageRotate / ImageRotateCW| & turn \\
\lstinline|ImageColorTint / Invert / Grayscale| & recolour every pixel \\
\lstinline|ImageColorBrightness / Contrast| & adjust \\
\lstinline|ImageColorReplace| & swap one exact colour for another \\
\lstinline|ImageDrawRectangle / Circle / Text| & \textbf{draw into the image} \\
\lstinline|ImageDraw| & paste one image into another \\
\lstinline|ImageFormat| & convert the pixel layout \\
\bottomrule
\end{tabular}
\end{center}

\gotcha{Almost all of these take a \lstinline|Image *| and modify it in place ---
\lstinline|ImageResize(&img, 128, 128)|, not \lstinline|img = ImageResize(...)|.
And they all work on the CPU, so they are far too slow to run every frame. Do
them once at load time, then upload.}

Use \lstinline|ImageCopy| before editing if you want to keep the original.

\section{Generating images from nothing}

\begin{lstlisting}
GenImageColor(w, h, colour);                      // a flat fill
GenImageGradientLinear(w, h, angle, a, b);        // a linear gradient
GenImageGradientRadial(w, h, density, a, b);      // a radial one: glows, vignettes
GenImageChecked(w, h, cw, ch, a, b);              // a checkerboard
GenImageWhiteNoise(w, h, factor);                 // static
GenImagePerlinNoise(w, h, offX, offY, scale);     // smooth noise: terrain, clouds
GenImageCellular(w, h, tileSize);                 // organic cell patterns
\end{lstlisting}

This is how you get a game looking like something before an artist exists.
\lstinline|GenImagePerlinNoise| feeding \lstinline|GenMeshHeightmap| from the
last chapter gives you an entire landscape in four lines.

\section{Render textures: drawing into a texture}

A \lstinline|RenderTexture2D| is a texture you can draw into, with its own
Begin/End pair:

\begin{lstlisting}
RenderTexture2D canvas = LoadRenderTexture(320, 240);     // once

// in the loop, BEFORE BeginDrawing:
BeginTextureMode(canvas);
    ClearBackground(BLACK);
    DrawCircle(160, 120, 40, ORANGE);
EndTextureMode();

BeginDrawing();
    DrawTexture(canvas.texture, 0, 0, WHITE);
EndDrawing();

UnloadRenderTexture(canvas);                              // once
\end{lstlisting}

This unlocks four things you cannot do otherwise:

\begin{enumerate}[leftmargin=*,itemsep=3pt]
  \item \textbf{Pixel-art at a fixed resolution.} Render your game at 320$\times$180
        and scale it up to the window. Every pixel stays square and crisp, at any
        window size.
  \item \textbf{A real minimap.} Draw the world from a top-down camera into a
        texture, then draw that texture in the corner.
  \item \textbf{Fullscreen shader effects.} Render the scene into a texture, then
        draw that texture through a shader. Bloom, colour grading, blur --- all
        of it works this way.
  \item \textbf{In-world screens.} Security monitors, mirrors, portals.
\end{enumerate}

\gotcha{A render texture comes out \textbf{upside down}, because OpenGL's Y axis
points the other way. The fix is a negative height in the source rectangle:

\lstinline|(Rectangle){ 0, 0, (float)rt.texture.width, -(float)rt.texture.height }|

Use \lstinline|DrawTexturePro| or \lstinline|DrawTextureRec| with that, and never
plain \lstinline|DrawTexture|, or your scene is mirrored vertically.}

\section{Texture filters and mipmaps}

\begin{lstlisting}
SetTextureFilter(tex, TEXTURE_FILTER_POINT);        // crisp: pixel art
SetTextureFilter(tex, TEXTURE_FILTER_BILINEAR);     // smooth: photos, 3D
GenTextureMipmaps(&tex);
SetTextureFilter(tex, TEXTURE_FILTER_TRILINEAR);    // smooth + mipmaps
\end{lstlisting}

\lstinline|POINT| is what makes pixel art look like pixel art; the default
bilinear filter will blur your beautiful 16$\times$16 sprite into mush the moment
you scale it.

Mipmaps are pre-shrunk copies of a texture. Without them, a wall texture seen
from far away shimmers as you walk, because each pixel is sampling a random spot
in the full-size image. Generate them for anything in a 3D world; skip them for
2D sprites drawn at their real size.

\section{Colour helpers}

\begin{center}
\begin{tabular}{@{}ll@{}}
\toprule
\lstinline|Fade(colour, 0.5f)| & same colour, half alpha \\
\lstinline|ColorTint(colour, tint)| & multiply two colours \\
\lstinline|ColorBrightness(colour, -0.4f)| & $-1$ black, $+1$ white \\
\lstinline|ColorContrast(colour, 0.5f)| & \\
\lstinline|ColorAlphaBlend(dst, src, tint)| & blend as if drawn on top \\
\lstinline|ColorLerp(a, b, t)| & blend two colours \\
\lstinline|ColorToHSV / ColorFromHSV| & hue, saturation, value \\
\lstinline|ColorToInt / GetColor| & to and from \lstinline|0xRRGGBBAA| \\
\bottomrule
\end{tabular}
\end{center}

HSV is the one worth knowing about. Shifting only the hue gives you a whole
palette of related colours --- team colours, damage states, day and night --- from
a single base:

\begin{lstlisting}
Vector3 hsv = ColorToHSV(base);
Color enemyTeam = ColorFromHSV(hsv.x + 180, hsv.y, hsv.z);   // opposite hue
\end{lstlisting}

And \lstinline|GetColor(0x1F2430FF)| lets you paste hex values straight from a
design tool, which is more pleasant than counting bytes.

\tryit{Module 6 shows all six image generators, a chain of CPU edits, the four
ways to draw a texture, a live render texture with its flipped source rectangle,
and the colour helpers side by side.}{./cheatsheet/build.sh 6}
